Publishing musical materials has always been a serious challenge. Niche subjects rarely pay for the kind of production staff a big house keeps on payroll, so the spreadsheet smiles on another beginner book with fresh graphics while your advanced idea starves. That is economics, not a verdict on your musicianship.

It used to cost blood and money: place every letter by hand, pay for a print run before you knew who would buy it, then eat the returns. Now a small crew---or one stubborn author---can match or beat old-line typography on a laptop, spend almost nothing beyond time, and reach readers through channels you already use. If you are not trying to squeeze margin from every page, giving the work away is often how niche teaching material actually gets read.

The habit of free, forkable, collaboratively improved resources did not start in conservatory offices. Programmers built it first: plain text, shared history, machines that rebuild the artifact whenever someone lands a fix. LilyPond treats a score as text; LaTeX treats the book the same way; GitHub is where the files live while many hands touch them. None of that writes your pedagogy for you. It removes the excuse that ``only a publisher could ship this.''

This \ProjectName tree is a working instance of that stack: prose and scores split into files you can argue about separately, a repeatable PDF build, bibliography, glossary, a front list of musical examples, and pre-styled pieces you can steal. The repository \texttt{README} is the blunt instrument: Docker, VS Code, Git, and your first \texttt{make pdf} in the order a nervous first-timer actually needs. My contact information is there if you want a human in the loop.
